% !TEX program = tectonic
% Long Xiao - AIGC / Streaming Audio-Video Generation / Video Agent Resume
\documentclass[10pt,a4paper]{ctexart}

\usepackage[left=1.28cm,right=1.28cm,top=0.84cm,bottom=0.32cm]{geometry}
\usepackage{xcolor}
\usepackage{graphicx}
\usepackage{tabularx}
\usepackage{array}
\usepackage{enumitem}
\usepackage{titlesec}
\usepackage{hyperref}
\usepackage{fontspec}
\usepackage{microtype}
\usepackage{ragged2e}

\providecommand{\AccentHex}{9A6A3A}
\providecommand{\LinkHex}{2F6B7C}
\definecolor{navy}{HTML}{22344A}
\definecolor{teal}{HTML}{\AccentHex}
\definecolor{linkblue}{HTML}{\LinkHex}
\definecolor{ink}{HTML}{293746}
\definecolor{muted}{HTML}{6B7784}
\definecolor{rulegray}{HTML}{D4DCE2}

\IfFontExistsTF{Times New Roman}{\setmainfont{Times New Roman}}{}
\IfFontExistsTF{Arial}{\setsansfont{Arial}}{}
\IfFontExistsTF{Songti SC}{\setCJKmainfont{Songti SC}}{\setCJKmainfont{FandolSong-Regular}}
\IfFontExistsTF{PingFang SC}{%
  \setCJKsansfont{PingFang SC}%
}{%
  \setCJKsansfont{FandolHei-Regular}%
}

\hypersetup{colorlinks=true,urlcolor=linkblue,linkcolor=linkblue}
\pagestyle{empty}
\setlength{\parindent}{0pt}
\setlength{\tabcolsep}{0pt}
\setlength{\parskip}{0pt}
\linespread{1.00}


\newcommand{\sectiontitle}[2]{%
  \vspace{0.22em}
  {\color{teal}\rule{2.2pt}{1.10em}}\hspace{0.48em}%
  {\sffamily\color{navy}\fontsize{12.7}{14}\selectfont\bfseries #1}%
  \hspace{0.62em}{\sffamily\color{muted}\fontsize{7.1}{8}\selectfont\MakeUppercase{#2}}\par
  \vspace{0.10em}{\color{rulegray}\rule{\linewidth}{0.55pt}}\vspace{0.10em}
}

\newcommand{\entryhead}[3]{%
  \begin{tabularx}{\linewidth}{@{}>{\sffamily\bfseries\color{navy}}X >{\sffamily\raggedleft\arraybackslash\color{muted}}p{3.05cm}@{}}
    #1\quad{\color{teal}\textbar}\quad #2 & #3
  \end{tabularx}\vspace{0.03em}
}

\newenvironment{tightitems}{%
  \begin{itemize}[leftmargin=1.20em,itemsep=0.02em,topsep=0.02em,parsep=0pt,partopsep=0pt,label=\textcolor{teal}{\raisebox{0.15ex}{\scriptsize\textbullet}}]
}{\end{itemize}\vspace{0.02em}}

\newcommand{\skillline}[2]{%
  \par\noindent\hangindent=1.72cm\hangafter=1%
  \makebox[1.62cm][l]{\sffamily\bfseries\color{navy}#1}#2\par\vspace{0.06em}%
}

\begin{document}
\color{ink}
\fontsize{10.3}{11.5}\selectfont


% Header
\begin{minipage}[t]{0.82\linewidth}
  \vspace{0pt}
  {\sffamily\fontsize{22}{24}\selectfont\bfseries\color{navy} 龙潇}%
  \hspace{0.55em}{\sffamily\fontsize{9}{10}\selectfont\color{muted}LONG XIAO}%
  \hspace{1.15em}\raisebox{-0.20cm}{\href{https://longxiao2001.github.io}{\includegraphics[width=1cm,height=1cm]{long_xiao_homepage_qr.png}}}\par
  \vspace{0.22em}
  {\sffamily\fontsize{12.4}{14}\selectfont\bfseries\color{navy}流式音视频联合生成 · 视频生成 · 视频生成 Agent}\par
  \vspace{0.36em}
  {\sffamily\color{muted}\fontsize{9.15}{10.9}\selectfont
  17383126759\quad\textbar\quad \href{mailto:1215497652@qq.com}{1215497652@qq.com}\quad\textbar\quad
  \href{https://longxiao2001.github.io}{longxiao2001.github.io}}
\end{minipage}\hfill
\begin{minipage}[t]{0.125\linewidth}
  \vspace{0pt}\raggedleft
  \fcolorbox{teal}{white}{\includegraphics[width=1.52cm,height=1.84cm,keepaspectratio]{long_xiao_photo.jpg}}
\end{minipage}


\vspace{-0.20em}
\sectiontitle{教育背景}{Education}

\begin{tabularx}{\linewidth}{@{}X >{\raggedleft\arraybackslash}p{3.2cm}@{}}
  {\sffamily\textbf{北京航空航天大学}}｜自动化科学与电气工程学院｜硕士｜{\sffamily\color{muted}专业排名 \textbf{\color{teal}前 10\%}} & {\sffamily\color{muted} 2024.09--至今}\\[0.09em]
  {\sffamily\textbf{东南大学}}｜自动化学院｜本科｜{\sffamily\color{muted}专业排名 \textbf{\color{teal}前 30\%}} & {\sffamily\color{muted} 2020.09--2024.06}
\end{tabularx}
\vspace{0.02em}

\sectiontitle{实习与研究经历}{Research Experience}
\entryhead{腾讯}{实时流式可交互音视频生成（LTX-2.3 22B）}{2026.06--2026.09}
\begin{tightitems}
  \item \textbf{基于LTX-2.3做实时交互式音视频生成}。把双向 LTX-2.3 改造为\textbf{块因果}模型。视频与音频的 latent 帧率不同（3 fps 与 25 fps），按 1 秒对齐设计 \textbf{AV 联合块}（视频 4+3k 帧、音频 26+25k 帧），并对\textbf{全部 6 路注意力}加因果掩码。\href{https://longxiao2001.github.io/StreamLTX/}{[Project Page]}
  \item 两阶段训练。\textbf{Teacher Forcing} 阶段使用真实长视频数据训练将双向扩散模型因果化；\textbf{Resample Forcing} 阶段用无梯度自回归 rollout 生成历史并在其上监督，缓解训练用真值加噪、推理用模型输出造成的误差累积。两阶段均用 concat 布局，两段共用一套 RoPE 索引并配稀疏因果掩码，单次前向得到 N 个块的监督。
  \item 长序列 KV 分 anchor、内容 memory、FIFO recent 三层。逐层按 QK 注意力质量与时间多样性打分做 Top-M 淘汰。
  \item 支持\textbf{流式生成过程中切换prompt}：设计\textbf{分层 Prompt 条件机制}，将整段文本条件解耦为描述场景、角色与风格的Global Prompt，以及承载阶段性台词与动作指令的Local Prompt。当前块只对当前片段prompt做交叉注意力计算。
  \item 推理\textbf{按去噪级切分多卡}，每卡常驻一级 KV，块间 NCCL P2P 传递latent 形成跨块流水；配合少步蒸馏，FA3，\textbf{4×H800、512×768} 下达到 \textbf{43 fps}（仅 DiT 稳态去噪吞吐），支持\textbf{分钟级}连续生成。
\end{tightitems}

\entryhead{阿里巴巴}{AIGC 视频生成算法｜时空局部重绘与表情编辑}{2025.10--2026.04}
\begin{tightitems}
  \item 基于 Wan2.2 做视频时空局部重绘，用\textbf{逐帧向量化 timestep 代替额外的掩码分支}：flow matching 中 σ 是逐位置的混合权重，取 0 保留原内容、取 1 完全重绘、取中间值做部分编辑，不改模型结构、不增加参数。
  \item 训练用\textbf{结构化重绘课程}（同步 / 随机异步 / 连续段重绘 / 边界平滑），比例随进度调整。硬阶跃的 timestep 会让 \textbf{adaLN 调制}在相邻帧突变、在编辑边界留下接缝，改用宽度 2 帧的线性过渡。LoRA rank 512，81 帧 480×832。
  \item Wan2.2 高/低噪专家以 σ=0.875 为界路由，与逐帧 timestep 冲突；按专家分别训练 LoRA，推理端按步路由、对条件帧按编辑强度一次性加噪，并沿缩放后的噪声轨迹去噪，可\textbf{指定任意时间段重绘而其余片段不变}。
  \item 搭建 UE Metahuman 表情参数提取与 3D 动画渲染管线，构建 2D 表情迁移 benchmark 与数据筛选流程。\href{https://youku-aigc.github.io/PerformRecast/}{[Project Page]}
\end{tightitems}

\entryhead{右脑科技}{2D Talking Avatar 与论文讲解视频 Agent}{2025.06--2025.10}
\begin{tightitems}
  \item 基于 VACE、LivePortrait 融合 Pose / Text / Audio 条件生成 2D 数字人，用分段生成与跨片段衔接扩展时长；用 CausVid 蒸馏把采样步数压到少步，配合显存动态加载，在\textbf{单张 RTX 4090} 上跑通VACE14B。
  \item 用 LangGraph 搭建\textbf{论文讲解视频 Agent}，链路为 MinerU MCP 解析 PDF、LLM 生成大纲与口播稿、渲染 HTML 幻灯片、TTS 配音、数字人口播、合成成片。用 \textbf{TTS 返回的字符级时间戳}确定每个分镜的时长，视觉素材按该时长生成，合成端只做小幅裁剪变速。
  \item 用 Send() 把分镜扇出为并行子图、以 operator.add 归并；每个子图带重试（报错信息回填进prompt）与确定性兜底工具，单个工具偶发失败时整条链路仍能输出。
\end{tightitems}

\entryhead{腾讯}{LLM 驱动的交互游戏 NPC}{2025.02--2025.06}
\begin{tightitems}
  \item 设计并集成 \textbf{ASR-LLM-TTS-UE} 端到端交互链路，完成模块封装、接口编排与推理流程联调，并在各环节间分配时延预算满足实时交互。
\end{tightitems}

\entryhead{本科毕业设计}{语音驱动的情感化 3D 数字人表情生成}{2024.01--2024.06}
\begin{tightitems}
  \item 冻结 \textbf{WavLM-Large} 编码语音，与 5 类情绪嵌入、前 8 帧种子表情共同作为条件，用 Transformer 去噪主干（MDM 结构）在 \textbf{185 维 Metahuman control rig} 上做扩散生成；预测 x0 并加入\textbf{速度损失}，抑制渲染后肉眼可见的抖动。
  \item 长音频按窗口滚动生成，上一段末尾 8 帧作下一段种子、重叠区线性混合，支持任意时长；导出 JSON 驱动 Metahuman 渲染成片。
\end{tightitems}

\sectiontitle{论文成果}{Publications}
\begin{enumerate}[leftmargin=1.5em,itemsep=0.04em,topsep=0.03em,parsep=0pt]
  \item J. Liang, B. Xiong, J. Tian, H. Li, \textbf{X. Long}, Y. Zheng, H. Fu. ``PerformRecast: Expression and Head Pose Disentanglement for Portrait Video Editing.'' \textit{CVPR 2026}. \href{https://youku-aigc.github.io/PerformRecast/}{[Project Page]}
  \item Z. Chen, \textbf{X. Long}, L. Zhang and W. Cai. ``Toward Diffusion-Based Deep Reinforcement Learning for Discrete Decision-Making: Methods and Evaluations.'' \textit{IEEE Transactions on Network Science and Engineering}.
\end{enumerate}

\sectiontitle{技术能力与综合亮点}{Skills \& Activities}
\skillline{技术方法}{Block-Causal Attention、Teacher / Diffusion / Resample Forcing、KV Cache 管理、少步蒸馏、时间步并行推理}
\skillline{模型框架}{LTX-2.3, Wan2.2, CogVideoX, MiniMax-H3, VACE, UE Metahuman, LangGraph}
\skillline{工程能力}{数据构建与筛选、训练实验设计、显存与推理优化、效果评估、系统集成服务部署}
% \skillline{音视频}{AV 联合块对齐、跨模态因果注意力（A2V / V2A）、latent 帧率对齐与 RoPE 设计、流式音频拼接}
% \skillline{生成模型}{Flow Matching / Diffusion、向量化时间步与局部重绘、LoRA、MoE 专家路由、adaLN 条件注入}
% \skillline{工程能力}{LangGraph 并行扇出与降级、MCP、FSDP / DeepSpeed、显存切片与权重动态加载、数据构建}
% \skillline{综合素质}{BEATBOX 协会会长；重庆市美术类联考 546 / 11900、北京电影学院第 32 名}

\end{document}
